% Part: computability
% Chapter: recursive-functions
% Section: examples

\documentclass[../../../include/open-logic-section]{subfiles}

\begin{document}

% SEGMENT OLP-0216-S01

\olfileid{cmp}{rec}{exa}
\olsection{முதன்மை மீள்வரையறைச் சார்புகளின் எடுத்துக்காட்டுகள்}


முதன்மை மீள்வரையறைச் சார்புகளுக்குச் சில எடுத்துக்காட்டுகள் ஏற்கெனவே
நம்மிடம் உள்ளன: கூட்டல் மற்றும் பெருக்கல் சார்புகள்~$\Add$ மற்றும்
$\Mult$. ஒன்றாதல் சார்பு $\fn{id}(x) = x$ ஒரு முதன்மை மீள்வரையறைச்
சார்பு; ஏனெனில் அது~$\Proj{1}{0}$ மட்டுமே. மாறிலிச் சார்புகள்
$\fn{const}_n(x) = n$ முதன்மை மீள்வரையறைச் சார்புகள்; ஏனெனில்
$\Zero$ மற்றும் $\Succ$ இலிருந்து அடுத்தடுத்த சேர்ப்புகளால் அவற்றை
வரையறுக்கலாம். முதன்மை மீள்வரையறைகளில் மாறிலிகளைப் பயன்படுத்த
விரும்பும்போது இது பயனுள்ளதாகும். எடுத்துக்காட்டாக, சார்பு
$f(x) = 2 \cdot x$ ஐ வரையறுக்க விரும்பினால்,
$\fn{const}_n(x)$ மற்றும் பெருக்கலிலிருந்து சேர்ப்பால்
$f(x) = \Mult(\fn{const}_2(x), \Proj{1}{0}(x))$ என அதைப் பெறலாம்.
இனி இத்தந்திரத்தைப் பயன்படுத்துவோம்.

\begin{prop}
  அடுக்கேற்றச் சார்பு $\fn{exp}(x, y) = x^y$ ஒரு முதன்மை
  மீள்வரையறைச் சார்பு.
\end{prop}

\begin{proof}
  $\fn{exp}$ ஐ முதன்மை மீள்வரையறை முறையில் பின்வருமாறு வரையறுக்கலாம்:
  \begin{align*}
    \fn{exp}(x, 0) & = 1\\
    \fn{exp}(x, y+1) & = \Mult(x, \fn{exp}(x,y)).
    \intertext{துல்லியமாகக் கூறினால், இது முதன்மை மீள்வரையறைச்
      சார்புகளிலிருந்து தரப்பட்ட மீள்வரையறை அல்ல. எனினும் முறைப்படி
      நம்மிடம் பின்வருவது உள்ளது:}
    \fn{exp}(x, 0) & = f(x)\\
    \fn{exp}(x, y+1) & = g(x, y, \fn{exp}(x,y)).
    \intertext{இங்கே}
    f(x) & = \Succ(\Zero(x)) = 1\\
    g(x, y, z) & = \Mult(\Proj{3}{0}(x, y, z), \Proj{3}{2}(x, y, z)) = x \cdot z
  \end{align*}
  எனவே $f$ மற்றும் $g$ முதன்மை மீள்வரையறைச் சார்புகளிலிருந்து
  சேர்ப்பால் வரையறுக்கப்படுகின்றன.
\end{proof}

\begin{prop}
  பின்வருமாறு வரையறுக்கப்படும் முன்னெண் சார்பு $\fn{pred}(y)$:
  \[
  \fn{pred}(y) = \begin{cases}
    0 & \text{$y=0$ எனில்}\\
    y-1 & \text{இல்லையெனில்}
  \end{cases}
  \]
  ஒரு முதன்மை மீள்வரையறைச் சார்பு.
\end{prop}

\begin{proof}
 பின்வருவனவற்றைக் கவனிக்க:
 \begin{align*}
   \fn{pred}(0) & = 0 \text{ மற்றும்}\\
   \fn{pred}(y+1) & = y.
 \end{align*}
 இது கிட்டத்தட்ட ஒரு முதன்மை மீள்வரையறை. துல்லியமாகக் கூறினால்,
 முதன்மை மீள்வரையறை வழி வரையறையின் முறை இதற்குப் பொருந்தவில்லை;
 ஏனெனில் அம்முறைக்கு குறைந்தபட்சம் ஒரு கூடுதல் உள்ளீடு~$x$
 தேவைப்படுகிறது. $\fn{pred}(y+1)$ இன் வரையறையில்
 $\fn{pred}(y)$ உண்மையில் பயன்படுத்தப்படாததும் வழக்கத்திற்கு மாறானது.
 ஆனால் முதலில் $\fn{pred}'(x, y)$ ஐப் பின்வருமாறு வரையறுக்கலாம்:
 \begin{align*}
   \fn{pred}'(x, 0) & = \Zero(x) = 0,\\
   \fn{pred}'(x, y+1) & = \Proj{3}{1}(x, y, \fn{pred'}(x, y)) = y.
 \end{align*}
பின்னர் அதிலிருந்து சேர்ப்பால் $\fn{pred}$ ஐ, எடுத்துக்காட்டாக,
$\fn{pred}(x) = \fn{pred}'(\Zero(x), \Proj{1}{0}(x))$ என
வரையறுக்கலாம்.
\end{proof}

\begin{prop}
  காரணியச் சார்பு $\fn{fac}(x) = \fact{x} = 1 \cdot 2 \cdot 3
  \cdot \dots \cdot x$ ஒரு முதன்மை மீள்வரையறைச் சார்பு.
\end{prop}

\begin{proof}
  தெளிவான முதன்மை மீள்வரையறை:
  \begin{align*}
    \fn{fac}(0) &= 1\\
    \fn{fac}(y+1) & = \fn{fac}(y) \cdot (y+1).
    \intertext{முறைப்படி, முதலில் ஓர் இருவிடச் சார்பு $h$ ஐ வரையறுக்க வேண்டும்:}
    h(x, 0) & = \fn{const}_1(x)\\
    h(x, y+1) & = g(x, y, h(x, y))
    \intertext{இங்கே $g(x, y, z) = \Mult(\Proj{3}{2}(x, y, z),
      \Succ(\Proj{3}{1}(x, y, z)))$; பின்னர் பின்வருமாறு வைக்க:}
    \fn{fac}(y) & = h(\Proj{1}{0}(y), \Proj{1}{0}(y)) = h(y,y).
  \end{align*}
  இனி சற்றுத் தளர்வாகச் செயல்பட்டு, சேர்ப்பு மற்றும் முதன்மை
  மீள்வரையறை வழியான முறைப்படியான வரையறைகளைத் தராமல் விடுவோம்.
\end{proof}

\begin{prop}
  பின்வருமாறு வரையறுக்கப்படும் குறைக்கப்பட்ட கழித்தல், $x \tsub y$:
  \[
  x \tsub y = \begin{cases}
    0 & \text{$x < y$ எனில்}\\
    x-y & \text{இல்லையெனில்}
  \end{cases}
  \]
  ஒரு முதன்மை மீள்வரையறைச் சார்பு.
\end{prop}

\begin{proof}
  நம்மிடம் பின்வருவன உள்ளன:
  \begin{align*}
    x \tsub 0 & = x\\
    x \tsub (y+1) & = \fn{pred}(x \tsub y) 
  \end{align*}
\end{proof}

\begin{prop}
 $x$ மற்றும் $y$ இடையிலான தொலைவு, $\left|x-y\right|$, ஒரு
 முதன்மை மீள்வரையறைச் சார்பு.
\end{prop}

\begin{proof}
  $\left| x-y \right| = (x \tsub y) + (y \tsub x)$ நம்மிடம்
  உள்ளது. எனவே முதன்மை மீள்வரையறைச் சார்புகளான $+$ மற்றும்
  $\tsub$ இலிருந்து சேர்ப்பால் தொலைவை வரையறுக்கலாம்.
\end{proof}

\begin{prop}
  $x$ மற்றும் $y$ இன் மீப்பெரு மதிப்பு, $\fn{max}(x,y)$, ஒரு
  முதன்மை மீள்வரையறைச் சார்பு.
\end{prop}

\begin{proof}
  $+$ மற்றும் $\tsub$ இலிருந்து சேர்ப்பால் $\fn{max}(x,y)$ ஐப்
  பின்வருமாறு வரையறுக்கலாம்:
  \[
  \fn{max}(x,y) \defis x + (y \tsub x).
  \]
  $x$ மீப்பெரு மதிப்பு, அதாவது $x \ge y$, எனில்
  $y \tsub x = 0$; எனவே $x + (y \tsub x) = x + 0 = x$.
  $y$ மீப்பெரு மதிப்பு எனில், $y \tsub x = y - x$; எனவே
  $x + (y \tsub x) = x + (y - x) = y$.
\end{proof}

\begin{prop}
  \ollabel{prop:min-pr}
  $x$ மற்றும் $y$ இன் மீச்சிறு மதிப்பு, $\fn{min}(x,y)$, ஒரு
  முதன்மை மீள்வரையறைச் சார்பு.
\end{prop}

\begin{proof}
  பயிற்சி.
\end{proof}

\begin{prob}
  \olref[cmp][rec][exa]{prop:min-pr} ஐ நிரூபிக்க.
\end{prob}

\begin{prob}
\[
f(x, y) =
2^{(2^{\iddots^{2^{x}}})}\raisebox{1ex}{\bigg\rbrace}
\raisebox{1ex}{\text {$y$ $2$'s}}\] ஒரு முதன்மை
மீள்வரையறைச் சார்பு என்று காட்டுக.
\end{prob}

\begin{prob}
முழுஎண் வகுத்தல் $d(x, y) = \lfloor x/y \rfloor$ (அதாவது, பதின்மப்
புள்ளிக்குப் பின்னுள்ள அனைத்தையும் புறக்கணிக்கும் வகுத்தல்) ஒரு
முதன்மை மீள்வரையறைச் சார்பு என்று காட்டுக. $y = 0$ எனில்,
$d(x, y) = 0$ என்று விதிக்கிறோம். முதன்மை மீள்வரையறை மற்றும்
சேர்ப்பைப் பயன்படுத்தி~$d$ இன் வெளிப்படையான வரையறையைத் தருக.
\end{prob}

\begin{prop}
முதன்மை மீள்வரையறைச் சார்புகளின் கணம் பின்வரும் இரு செயல்களின்
கீழ் மூடியது:
\begin{enumerate}
\item இறுதிக்குட்பட்ட கூட்டுத்தொகைகள்: $f(\vec x, z)$ ஒரு முதன்மை
மீள்வரையறைச் சார்பு எனில், பின்வரும் சார்பும் அவ்வாறே அமையும்:
\[
g(\vec x, y) \defis \sum_{z = 0}^y f(\vec x, z).
\]
\item இறுதிக்குட்பட்ட பெருக்குத்தொகைகள்: $f(\vec x, z)$ ஒரு முதன்மை
மீள்வரையறைச் சார்பு எனில், பின்வரும் சார்பும் அவ்வாறே அமையும்:
\[
h(\vec x, y) \defis \prod_{z = 0}^y f(\vec x, z).
\]
\end{enumerate}
\end{prop}

\begin{proof}
எடுத்துக்காட்டாக, இறுதிக்குட்பட்ட கூட்டுத்தொகைகள் பின்வரும்
சமன்பாடுகளால் மீள்வரையறை முறையில் வரையறுக்கப்படுகின்றன:
\begin{align*}
  g(\vec x, 0) & = f(\vec x, 0)\\
  g(\vec x, y+1) & = g(\vec x, y) + f(\vec x, y+1).
\end{align*}
\end{proof}


\end{document}
